
\mectable%
  {Individual Outcomes: Intensity$_{1998}$ $\times$ Post, All Municipalities (1992--1998)}%
  {tab:t1_ind_1992_1998_all}%
  {tables/1998/T1_ind_enigh_1992_1998_all}%
  {This table shows the results obtained from estimating a difference-in-differences model using individual labor-related outcomes. It includes municipality and year fixed effects. Standard errors are clustered at the level of treatment: municipality. $Intensity_{1998}$ captures the intervention effect due to the intensity of the
treatment on individual labor-related outcomes, i.e., the effect when 0\% of households in a municipality are Progresa beneficiaries goes up to 100\% after 1998. The period of analysis is from 1992 to 1998, using survey waves available every two years. It includes all municipalities. All regressions are weighted using the survey weights, and excludes upper 1\% outliers. All monetary values are inflation-adjusted up to October 2025. ***p$<0.01$, **p$<0.05$, *p$<0.1$}%

\mectable%
  {Household Outcomes: Intensity$_{1998}$ $\times$ Post, All Municipalities (1992--1998)}%
  {tab:t2_hh_1992_1998_all}%
  {tables/1998/T2_hh_enigh_1992_1998_all}%
  {This table shows the results obtained from estimating a difference-in-differences model using household labor-related outcomes. It includes municipality and year fixed effects. Standard errors are clustered at the level of treatment: municipality. $Intensity_{1998}$ captures the intervention effect due to the intensity of the treatment on household labor outcomes, i.e., the effect when 0\% of households in a municipality are Progresa beneficiaries goes up to 100\% after 1998. The period of analysis is from 1992 to 1998, using survey waves available every two years. It includes all municipalities. All regressions are weighted using the survey weights, and excludes upper 1\% outliers. All monetary values are inflation-adjusted up to October 2025. ***p$<0.01$, **p$<0.05$, *p$<0.1$}%


\mectable%
  {Food Expenditure Outcomes: Intensity$_{1998}$ $\times$ Post, All Municipalities (1992--1998)}%
  {tab:t3_food_1992_1998_all}%
  {tables/1998/T3_food_enigh_1992_1998_all}%
  {This table shows the results obtained from estimating a difference-in-differences model using household food expenditure outcomes. It includes municipality and year fixed effects. Standard errors are clustered at the level of treatment: municipality. $Intensity_{1998}$ captures the intervention effect due to the intensity of the treatment on household food expenditure, i.e., the effect when 0\% of households in a municipality are Progresa beneficiaries goes up to 100\% after 1998. The period of analysis is from 1992 to 1998, using survey waves available every two years. It includes all municipalities. All regressions are weighted using the survey weights, and excludes upper 1\% outliers. All monetary values are inflation-adjusted up to October 2025. ***p$<0.01$, **p$<0.05$, *p$<0.1$}%


\mectable%
  {Health Expenditure Outcomes: Intensity$_{1998}$ $\times$ Post, All Municipalities (1992--1998)}%
  {tab:t4_health_1992_1998_all}%
  {tables/1998/T4_health_enigh_1992_1998_all}%
  {This table shows the results obtained from estimating a difference-in-differences model using household health expenditure outcomes. It includes municipality and year fixed effects. Standard errors are clustered at the level of treatment: municipality. $Intensity_{1998}$ captures the intervention effect due to the intensity of the treatment on household health expenditure, i.e., the effect when 0\% of households in a municipality are Progresa beneficiaries goes up to 100\% after 1998. The period of analysis is from 1992 to 1998, using survey waves available every two years. It includes all municipalities. All regressions are weighted using the survey weights, and excludes upper 1\% outliers. All monetary values are inflation-adjusted up to October 2025. ***p$<0.01$, **p$<0.05$, *p$<0.1$}%




\mectable%
  {Individual Outcomes: Intensity$_{2000}$ $\times$ Post, All Municipalities (1992--2000)}%
  {tab:t1_ind_1992_2000_all}%
  {tables/2000/T1_ind_enigh_1992_2000_all}%
  {This table shows the results obtained from estimating a difference-in-differences model using individual labor-related outcomes. It includes municipality and year fixed effects. Standard errors are clustered at the level of treatment: municipality. $Intensity_{2000}$ captures the intervention effect due to the intensity of the treatment on household health expenditure, i.e., the effect when 0\% of households in a municipality are Progresa beneficiaries goes up to 100\% after 2000. The period of analysis is from 1992 to 2000, using survey waves available every two years, and excludes 1998. It includes all municipalities. All regressions are weighted using the survey weights, and excludes upper 1\% outliers. All monetary values are inflation-adjusted up to October 2025. ***p$<0.01$, **p$<0.05$, *p$<0.1$}%
  
\mectable%
  {Household Outcomes: Intensity$_{2000}$ $\times$ Post, All Municipalities (1992--2000)}%
  {tab:t2_hh_1992_2000_all}%
  {tables/2000/T2_hh_enigh_1992_2000_all}%
  {This table shows the results obtained from estimating a difference-in-differences model using household labor-related outcomes. It includes municipality and year fixed effects. Standard errors are clustered at the level of treatment: municipality. $Intensity_{2000}$ captures the intervention effect due to the intensity of the treatment on household health expenditure, i.e., the effect when 0\% of households in a municipality are Progresa beneficiaries goes up to 100\% after 2000. The period of analysis is from 1992 to 2000, using survey waves available every two years, and excludes 1998. It includes all municipalities. All regressions are weighted using the survey weights, and excludes upper 1\% outliers. All monetary values are inflation-adjusted up to October 2025. ***p$<0.01$, **p$<0.05$, *p$<0.1$}%

\mectable%
  {Food Expenditure Outcomes:Intensity$_{2000}$ $\times$ Post, All Municipalities (1992--2000)}%
  {tab:t3_food_1992_2000_all}%
  {tables/2000/T3_food_enigh_1992_2000_all}%
  {This table shows the results obtained from estimating a difference-in-differences model using household food expenditure outcomes. It includes municipality and year fixed effects. Standard errors are clustered at the level of treatment: municipality. $Intensity_{2000}$ captures the intervention effect due to the intensity of the treatment on household health expenditure, i.e., the effect when 0\% of households in a municipality are Progresa beneficiaries goes up to 100\% after 2000. The period of analysis is from 1992 to 2000, using survey waves available every two years, and excludes 1998. It includes all municipalities. All regressions are weighted using the survey weights, and excludes upper 1\% outliers. All monetary values are inflation-adjusted up to October 2025. ***p$<0.01$, **p$<0.05$, *p$<0.1$}%

\mectable%
  {Health Expenditure Outcomes: Intensity$_{2000}$ $\times$ Post, All Municipalities (1992--2000)}%
  {tab:t4_health_1992_2000_all}%
  {tables/2000/T4_health_enigh_1992_2000_all}%
  {This table shows the results obtained from estimating a difference-in-differences model using household health expenditure outcomes. It includes municipality and year fixed effects. Standard errors are clustered at the level of treatment: municipality. $Intensity_{2000}$ captures the intervention effect due to the intensity of the treatment on household health expenditure, i.e., the effect when 0\% of households in a municipality are Progresa beneficiaries goes up to 100\% after 2000. The period of analysis is from 1992 to 2000, using survey waves available every two years, and excludes 1998. It includes all municipalities. All regressions are weighted using the survey weights, and excludes upper 1\% outliers. All monetary values are inflation-adjusted up to October 2025. ***p$<0.01$, **p$<0.05$, *p$<0.1$}%
